\section{Introduction}

The economic question is not only whether artificial intelligence can perform
tasks currently done by workers. It is also whether AI can perform the research
that produces better AI. Once the same technology enters both current production
and the improvement of future technology, a familiar substitution problem becomes
a dynamic feedback problem. Better AI lowers the compute cost of AI services,
raises output and the resources available for research, and makes research compute
itself more productive. Whether that loop settles into balanced growth or instead
accelerates depends on where human labor remains essential.

This paper studies that feedback in a general-equilibrium growth model. Raw
compute has two uses. Inference compute is transformed by current AI capability
into services used in final production. Research compute is transformed by the
same capability into automated research services. Human researchers and automated
research services are combined with a CES technology to improve AI capability.
An integrated monopolist supplies AI services and chooses its research inputs.
Households own capital and the developer, supply labor, and choose consumption.
Population grows exogenously.

The distinction between raw compute and the services produced by compute is
central. A unit of compute is not a unit of intelligence. Current capability
determines what a given unit of hardware, energy, and datacenter time can produce.
We therefore write production services as the product of capability and inference
compute, and automated research services as the product of capability and research
compute. This symmetric treatment corresponds closely to the distinction between
inference and training compute in \citet{korinekmckelvey2026}. It also removes a
potential inconsistency in models that let capability augment compute in final
production but not in AI research.

The model deliberately excludes a separate standing-on-shoulders term from its
baseline research technology. Better AI affects research only by making research
compute more productive. This choice prevents the model from counting the same
technological feedback twice and makes the source of acceleration transparent. A
direct knowledge-spillover term can be added later, but it is not needed for the
mechanism studied here.

The analysis separates three regimes according to the elasticity of substitution
between labor and AI services in final production. When the elasticity is below
one, the monopoly supplies a finite ratio of AI services to labor as capability
becomes abundant. The economy then approaches a production technology in capital
and labor alone, so output per capita cannot grow forever along a conventional
balanced path. At unit elasticity, the paper derives the growth rates and resource
shares of balanced-growth candidates. Automating research tightens the relevant
stability condition because capability now raises both inference services and
automated research services. When the elasticity exceeds one, unbounded capability
raises output relative to capital and hence the return to capital without bound.
A finite-rate steady state is impossible.

For the gross-substitutes regime, the paper goes beyond this nonexistence result.
It characterizes an AI-dominated limiting equilibrium in which output, capital,
consumption, and capability accelerate. The result is conditional on the path
remaining interior and on the monopoly solution satisfying its second-order
condition. Within the parameter region for which that limit is established,
labor's share tends to zero, while the real wage rises below an explicit
substitution threshold and falls above it. This distinction
between labor's share and the wage is economically important: a factor may receive
a vanishing fraction of an economy that is expanding even faster.

Perfect-foresight simulations solve the household Euler equation, the developer's
research first-order conditions and costate equation, the monopoly pricing
condition, both market-clearing conditions, and the two laws of motion. The
simulations are used to illustrate propositions, not to replace them. The reported
paths are rejected unless their resource, optimality, and dynamic residuals are
small.

The paper contributes in three ways. First, it provides a symmetric accounting of
inference and research compute. Second, it shows exactly how this accounting
changes the balanced-growth denominators: capability enters automated research
through the product of capability and research compute. Third, it separates the
long-run behavior of wages from the long-run behavior of labor's income share on
an accelerating equilibrium path. These results connect the literature on
research-based growth, automation, factor substitution, and economic singularities
within one tractable model.
